\chapter{Équations, inéquations et systèmes du second degré}
\textit{Le second degré est l'un des outils les plus puissants du programme de
Première C. Dès qu'une vitesse, une aire ou une trajectoire fait intervenir un
carré, on retombe sur une expression de la forme $ax^2+bx+c$. Apprendre à la
factoriser, à en lire le signe et à résoudre les équations associées, c'est se
doter d'un réflexe que l'on retrouvera en géométrie, en physique et jusqu'au
calcul intégral. Tout repose sur une seule idée : la \emph{forme canonique}.}

\section{Le trinôme du second degré}

\begin{definition}[Trinôme du second degré]
On appelle \textbf{trinôme du second degré} toute expression de la forme
\[ P(x)=ax^2+bx+c, \qquad a,b,c\in\R,\ a\neq 0. \]
Le réel $a$ est le \textbf{coefficient dominant} ; la condition $a\neq 0$ est
essentielle (sinon l'expression est du premier degré). L'\textbf{équation du second
degré} associée est $ax^2+bx+c=0$.
\end{definition}

\begin{remarque}
Si $a=0$, l'équation $bx+c=0$ est du premier degré et se résout immédiatement.
Tout ce qui suit suppose donc $a\neq 0$ : on parle bien du \emph{second} degré.
\end{remarque}

\begin{theoreme}[Forme canonique]
Pour tout réel $x$ et tout trinôme $P(x)=ax^2+bx+c$ avec $a\neq 0$,
\[ ax^2+bx+c = a\!\left[\left(x+\frac{b}{2a}\right)^{\!2}-\frac{b^2-4ac}{4a^2}\right]. \]
\end{theoreme}

\begin{aretenir}[Démonstration de la forme canonique]
On factorise $a$, puis on reconnaît le début d'un carré :
\[ ax^2+bx+c = a\!\left(x^2+\frac{b}{a}x\right)+c
   = a\!\left[\left(x+\frac{b}{2a}\right)^{\!2}-\frac{b^2}{4a^2}\right]+c. \]
En remettant $c=a\cdot\dfrac{4ac}{4a^2}$ à l'intérieur du crochet, il vient
\[ ax^2+bx+c = a\!\left[\left(x+\frac{b}{2a}\right)^{\!2}-\frac{b^2-4ac}{4a^2}\right]. \]
Toute la résolution du second degré découle de cette identité.
\end{aretenir}

\begin{exemple}
Mettons $P(x)=2x^2-12x+10$ sous forme canonique. On factorise $2$ :
\[ P(x)=2(x^2-6x+5)=2\big[(x-3)^2-9+5\big]=2(x-3)^2-8. \]
On lit immédiatement que $P$ admet un minimum égal à $-8$, atteint en $x=3$.
\end{exemple}

\section{Le discriminant et les racines}

La quantité $b^2-4ac$ apparue dans la forme canonique gouverne entièrement le
nombre de solutions de l'équation $ax^2+bx+c=0$.

\begin{definition}[Discriminant]
On appelle \textbf{discriminant} du trinôme $ax^2+bx+c$ le réel
\[ \Delta = b^2-4ac. \]
\end{definition}

\begin{theoreme}[Nombre et valeurs des racines]
Soit $ax^2+bx+c=0$ avec $a\neq 0$ et $\Delta=b^2-4ac$.
\begin{itemize}
  \item Si $\Delta>0$ : l'équation a \textbf{deux racines réelles distinctes}
  \[ x_1=\frac{-b-\sqrt{\Delta}}{2a}, \qquad x_2=\frac{-b+\sqrt{\Delta}}{2a}. \]
  \item Si $\Delta=0$ : l'équation a \textbf{une racine double}
  \[ x_0=-\frac{b}{2a}. \]
  \item Si $\Delta<0$ : l'équation n'a \textbf{aucune racine réelle}.
\end{itemize}
\end{theoreme}

\begin{aretenir}[Pourquoi $\Delta$ décide de tout]
À partir de la forme canonique, $ax^2+bx+c=0$ équivaut (en divisant par $a\neq 0$) à
\[ \left(x+\frac{b}{2a}\right)^{\!2}=\frac{\Delta}{4a^2}. \]
Un carré ne peut égaler que des quantités positives ou nulles : si $\Delta<0$, le
second membre est négatif, donc aucune solution ; si $\Delta=0$, on obtient une
unique valeur ; si $\Delta>0$, on extrait deux racines carrées opposées, d'où les
deux solutions ci-dessus.
\end{aretenir}

\begin{illus}
\begin{tikzpicture}[scale=0.78]
  \begin{scope}
    \draw[->,onbline] (-1.6,0)--(2.6,0) node[right,onbmuted]{$x$};
    \draw[->,onbline] (0,-1.4)--(0,2.6);
    \draw[thick,onbo,smooth,samples=80,domain=-0.7:1.9] plot (\x,{2*(\x-0.6)*(\x-0.6)-1});
    \filldraw[onbink] (-0.107,0) circle (1.4pt);
    \filldraw[onbink] (1.307,0) circle (1.4pt);
    \node[onbmuted,below] at (0.6,-1.7) {$\Delta>0$ : deux racines};
  \end{scope}
  \begin{scope}[xshift=5cm]
    \draw[->,onbline] (-1.6,0)--(2.6,0) node[right,onbmuted]{$x$};
    \draw[->,onbline] (0,-0.6)--(0,2.6);
    \draw[thick,onbo,smooth,samples=80,domain=-0.25:1.45] plot (\x,{2.2*(\x-0.6)*(\x-0.6)});
    \filldraw[onbink] (0.6,0) circle (1.4pt);
    \node[onbmuted,below] at (0.6,-0.9) {$\Delta=0$ : racine double};
  \end{scope}
  \begin{scope}[xshift=10cm]
    \draw[->,onbline] (-1.6,0)--(2.6,0) node[right,onbmuted]{$x$};
    \draw[->,onbline] (0,-0.6)--(0,2.6);
    \draw[thick,onbo,smooth,samples=80,domain=-0.15:1.35] plot (\x,{2.2*(\x-0.6)*(\x-0.6)+0.7});
    \node[onbmuted,below] at (0.6,-0.9) {$\Delta<0$ : aucune racine};
  \end{scope}
\end{tikzpicture}
\fig{Figure — la parabole coupe l'axe $(Ox)$ en deux points, le touche, ou l'évite, selon le signe de $\Delta$.}
\end{illus}

\begin{exemple}
Résolvons $x^2-5x+6=0$. Ici $a=1$, $b=-5$, $c=6$, donc
$\Delta=(-5)^2-4\times 1\times 6=25-24=1>0$. D'où
\[ x_1=\frac{5-1}{2}=2, \qquad x_2=\frac{5+1}{2}=3. \]
L'ensemble des solutions est $S=\{2\,;\,3\}$.
\end{exemple}

\begin{remarque}
Lorsque $b$ est pair, on peut alléger les calculs avec le \textbf{discriminant
réduit} $\Delta'=\left(\frac{b}{2}\right)^2-ac$ : si $\Delta'>0$, les racines sont
$x=\dfrac{-\frac{b}{2}\pm\sqrt{\Delta'}}{a}$. C'est commode mais facultatif.
\end{remarque}

\section{Somme, produit et factorisation}

\begin{theoreme}[Somme et produit des racines]
Lorsque le trinôme $ax^2+bx+c$ admet deux racines $x_1$ et $x_2$ (distinctes si
$\Delta>0$, confondues si $\Delta=0$), alors
\[ S=x_1+x_2=-\frac{b}{a}, \qquad P=x_1\,x_2=\frac{c}{a}. \]
\end{theoreme}

\begin{aretenir}[Vérification rapide]
En additionnant les deux racines, $\sqrt{\Delta}$ se simplifie :
\[ x_1+x_2=\frac{-b-\sqrt\Delta}{2a}+\frac{-b+\sqrt\Delta}{2a}=\frac{-2b}{2a}=-\frac{b}{a}. \]
En les multipliant, on reconnaît une identité remarquable :
\[ x_1\,x_2=\frac{(-b)^2-(\sqrt\Delta)^2}{4a^2}=\frac{b^2-\Delta}{4a^2}=\frac{4ac}{4a^2}=\frac{c}{a}. \]
Ces relations permettent souvent de \emph{deviner} les racines sans calculer $\Delta$.
\end{aretenir}

\begin{theoreme}[Factorisation du trinôme]
\begin{itemize}
  \item Si $\Delta>0$, de racines $x_1,x_2$ : \quad $ax^2+bx+c=a(x-x_1)(x-x_2)$.
  \item Si $\Delta=0$, de racine double $x_0$ : \quad $ax^2+bx+c=a(x-x_0)^2$.
  \item Si $\Delta<0$ : le trinôme est \textbf{irréductible} dans $\R$ (pas de
        factorisation en facteurs du premier degré).
\end{itemize}
\end{theoreme}

\begin{exemple}
Le trinôme $x^2-5x+6$ a pour racines $2$ et $3$ (exemple précédent), donc se
factorise en $x^2-5x+6=(x-2)(x-3)$. Vérification immédiate : $S=2+3=5=-\frac{-5}{1}$
et $P=2\times 3=6=\frac{6}{1}$.
\end{exemple}

\begin{remarque}
Si $a+b+c=0$, alors $x=1$ est racine évidente et l'autre racine vaut $\frac{c}{a}$
(car $P=\frac{c}{a}$). Par exemple $2x^2-3x+1=0$ : $2-3+1=0$, donc $x=1$ et
$x'=\frac{1}{2}$. Ce réflexe fait gagner un temps précieux en composition.
\end{remarque}

\section{Signe du trinôme}

\begin{theoreme}[Signe du trinôme du second degré]
Le trinôme $P(x)=ax^2+bx+c$ ($a\neq 0$) est :
\begin{itemize}
  \item \textbf{toujours du signe de $a$} si $\Delta<0$ (aucune racine) ;
  \item du signe de $a$ \textbf{sauf en $x_0$} où il s'annule, si $\Delta=0$ ;
  \item du signe de $a$ \textbf{à l'extérieur} des racines, et du signe de $-a$
        (signe contraire) \textbf{entre les racines}, si $\Delta>0$.
\end{itemize}
On résume cela par la formule : \emph{« le trinôme est du signe de $a$, sauf
entre les racines »}.
\end{theoreme}

\begin{illus}
\begin{tikzpicture}[scale=0.85]
  \draw[->,onbline] (-2.6,0)--(2.8,0) node[right,onbmuted]{$x$};
  \draw[->,onbline] (0,-1.6)--(0,2.2) node[above,onbmuted]{$y$};
  \draw[thick,onbo,smooth,samples=90,domain=-1.65:1.65] plot (\x,{1.2*(\x-1)*(\x+1)});
  \filldraw[onbink] (-1,0) circle (1.5pt) node[below left,onbink]{$x_1$};
  \filldraw[onbink] (1,0) circle (1.5pt) node[below right,onbink]{$x_2$};
  \node[onbgreen] at (-1.7,1.4) {$+$};
  \node[onbgreen] at (1.7,1.4) {$+$};
  \node[onbo] at (0,-1.3) {$-$};
\end{tikzpicture}
\fig{Figure — cas $a>0$ et $\Delta>0$ : le trinôme est positif à l'extérieur de $[x_1\,;\,x_2]$ et négatif entre les racines.}
\end{illus}

\begin{exemple}
Étudions le signe de $P(x)=-x^2+x+6$. On a $a=-1<0$ et
$\Delta=1+24=25>0$, d'où les racines $x_1=\dfrac{-1-5}{-2}=3$ et
$x_2=\dfrac{-1+5}{-2}=-2$. Le tableau de signes :
\begin{center}\renewcommand{\arraystretch}{1.3}
\begin{tabular}{c|ccccc}
$x$ & $-\infty$ & $-2$ & & $3$ & $+\infty$ \\
\hline
$P(x)$ & & $-$ & $\;0\;+\;0\;$ & & $-$ \\
\end{tabular}
\end{center}
Comme $a<0$, $P$ est négatif à l'extérieur de $[-2\,;\,3]$ et positif entre $-2$ et $3$.
\end{exemple}

\section{Inéquations du second degré}

\begin{aretenir}[Méthode générale]
Pour résoudre une inéquation du second degré (par exemple $ax^2+bx+c\ge 0$) :
\begin{enumerate}
  \item ramener tout dans un seul membre pour obtenir $\,ax^2+bx+c \;\square\; 0$ ;
  \item calculer $\Delta$ et, s'il y a lieu, les racines ;
  \item dresser le \textbf{tableau de signes} du trinôme ;
  \item lire l'ensemble solution selon le sens de l'inégalité (attention aux bornes :
        $\ge$ ou $\le$ \emph{incluent} les racines, $>$ ou $<$ les \emph{excluent}).
\end{enumerate}
\end{aretenir}

\begin{exemple}
Résolvons $x^2-x-6<0$. Ici $\Delta=1+24=25$, racines $x_1=-2$ et $x_2=3$, avec
$a=1>0$ : le trinôme est positif à l'extérieur et négatif entre les racines. On
cherche $P(x)<0$, donc
\[ S=\,]-2\,;\,3[. \]
Les bornes sont exclues car l'inégalité est stricte.
\end{exemple}

\begin{remarque}
Quand $\Delta<0$, le trinôme garde un signe constant (celui de $a$). Ainsi
$x^2+x+1>0$ est vraie pour \emph{tout} $x\in\R$ (car $a>0$ et $\Delta=-3<0$), tandis
que $x^2+x+1<0$ n'a \emph{aucune} solution : $S=\varnothing$.
\end{remarque}

\section{Équations se ramenant au second degré}

Certaines équations de degré supérieur deviennent du second degré après un
\textbf{changement d'inconnue} bien choisi.

\begin{definition}[Équation bicarrée]
On appelle \textbf{équation bicarrée} toute équation de la forme
\[ ax^4+bx^2+c=0, \qquad a\neq 0. \]
En posant $X=x^2$ (avec la contrainte $X\ge 0$), elle se ramène à l'équation du
second degré $aX^2+bX+c=0$.
\end{definition}

\begin{aretenir}[Marche à suivre pour une bicarrée]
\begin{enumerate}
  \item poser $X=x^2$, avec $X\ge 0$ ;
  \item résoudre $aX^2+bX+c=0$ en $X$ ;
  \item ne garder que les solutions $X\ge 0$ ;
  \item revenir à $x$ : chaque $X>0$ donne $x=\pm\sqrt{X}$ ; $X=0$ donne $x=0$.
\end{enumerate}
\end{aretenir}

\begin{exemple}
Résolvons $x^4-5x^2+4=0$. On pose $X=x^2$ : $X^2-5X+4=0$, de racines $X=1$ et $X=4$
(somme $5$, produit $4$), toutes deux positives. On revient à $x$ :
\[ x^2=1\Rightarrow x=\pm 1, \qquad x^2=4\Rightarrow x=\pm 2. \]
D'où $S=\{-2\,;\,-1\,;\,1\,;\,2\}$.
\end{exemple}

\begin{remarque}
Le même réflexe s'applique à d'autres formes : poser $X=\sqrt{x}$ dans
$x-3\sqrt{x}+2=0$, ou $X=\dfrac{1}{x}$ dans une équation où apparaissent $x^2$ et
$\frac{1}{x^2}$. L'idée est toujours de \emph{nommer} le motif qui se répète pour
faire surgir un second degré.
\end{remarque}

\section{Systèmes du second degré}

\begin{theoreme}[Système somme--produit]
Deux réels $u$ et $v$ de somme $S$ et de produit $P$ sont les \textbf{racines} du
trinôme
\[ t^2-St+P=0. \]
Le système $\begin{cases} u+v=S\\ uv=P \end{cases}$ admet donc des solutions réelles
si et seulement si $S^2-4P\ge 0$ ; les valeurs de $u$ et $v$ sont alors les deux
racines de ce trinôme.
\end{theoreme}

\begin{exemple}
Cherchons deux nombres dont la somme vaut $7$ et le produit $12$. Ils sont racines
de $t^2-7t+12=0$, soit $\Delta=49-48=1$, racines $3$ et $4$. Les deux nombres sont
donc $3$ et $4$.
\end{exemple}

\begin{aretenir}[Système d'une équation de degré 1 et d'une de degré 2]
Pour un système du type $\begin{cases} x+y=s\\ x^2+y^2=k \end{cases}$ ou avec une
droite et une parabole, on procède par \textbf{substitution} : tirer une inconnue de
l'équation du premier degré, la reporter dans la seconde, puis résoudre le second
degré obtenu.
\end{aretenir}

\begin{exemple}
Résolvons $\begin{cases} x+y=5\\ xy=6 \end{cases}$. C'est un système
somme--produit : $x$ et $y$ sont racines de $t^2-5t+6=0$, soit $t=2$ ou $t=3$.
Les couples solutions sont $(2\,;\,3)$ et $(3\,;\,2)$.
\end{exemple}

% ════════════════════════════════════════════════════════════
\section{L'essentiel à retenir}

\begin{synthese}
\textbf{Forme canonique.} $ax^2+bx+c=a\!\left[\left(x+\frac{b}{2a}\right)^2-\frac{\Delta}{4a^2}\right]$
avec $\Delta=b^2-4ac$.

\smallskip
\textbf{Racines selon $\Delta$.}
$\Delta>0$ : deux racines $x_{1,2}=\dfrac{-b\pm\sqrt\Delta}{2a}$ ;
$\Delta=0$ : racine double $x_0=-\dfrac{b}{2a}$ ;
$\Delta<0$ : aucune racine réelle.

\smallskip
\textbf{Somme et produit.} $S=x_1+x_2=-\dfrac{b}{a}$, \ $P=x_1x_2=\dfrac{c}{a}$.
Astuce : si $a+b+c=0$ alors $x=1$ est racine.

\smallskip
\textbf{Factorisation.} $\Delta>0$ : $a(x-x_1)(x-x_2)$ ; \ $\Delta=0$ : $a(x-x_0)^2$ ;
\ $\Delta<0$ : irréductible dans $\R$.

\smallskip
\textbf{Signe.} Le trinôme est \emph{du signe de $a$, sauf entre les racines}
(quand $\Delta>0$) ; du signe de $a$ partout si $\Delta\le 0$.

\smallskip
\textbf{Bicarrée \& cie.} On pose $X=x^2$ (ou autre motif), on résout en $X\ge 0$,
puis on revient à $x$.

\smallskip
\textbf{Somme--produit.} $u,v$ de somme $S$ et produit $P$ sont racines de
$t^2-St+P=0$.
\end{synthese}

% ════════════════════════════════════════════════════════════
\section{Méthodes \& savoir-faire}

\methode{Mettre un trinôme sous forme canonique}
Factoriser $a$, compléter le carré à partir de $x^2+\frac{b}{a}x$, puis réintégrer
le terme constant. On lit alors directement le sommet de la parabole.
\begin{exemple}
$P(x)=x^2-4x+1=(x-2)^2-4+1=(x-2)^2-3$ : sommet $(2\,;\,-3)$, minimum $-3$.
\end{exemple}

\methode{Résoudre une équation du second degré}
Identifier $a,b,c$, calculer $\Delta=b^2-4ac$, conclure selon son signe et, si
$\Delta\ge 0$, appliquer $x=\dfrac{-b\pm\sqrt\Delta}{2a}$.
\begin{exemple}
$3x^2-2x-1=0$ : $\Delta=4+12=16$, $\sqrt\Delta=4$, d'où $x=\dfrac{2\pm 4}{6}$, soit
$x=1$ ou $x=-\dfrac{1}{3}$. On vérifie avec $a+b+c=3-2-1=0$ : $x=1$ est bien racine.
\end{exemple}

\methode{Trouver les racines par la somme et le produit}
Quand les racines sont « simples », chercher deux nombres de somme $-\frac{b}{a}$
et de produit $\frac{c}{a}$ — souvent plus rapide que le discriminant.
\begin{exemple}
$x^2-7x+10=0$ : on cherche deux nombres de somme $7$ et produit $10$ : ce sont $2$
et $5$. Donc $S=\{2\,;\,5\}$.
\end{exemple}

\methode{Résoudre une inéquation du second degré}
Tout ramener d'un côté, calculer $\Delta$ et les racines, dresser le tableau de
signes, puis lire l'ensemble solution selon le sens (strict ou large) de l'inégalité.
\begin{exemple}
$2x^2-3x-2\ge 0$ : $\Delta=9+16=25$, racines $-\frac{1}{2}$ et $2$, $a>0$. Le trinôme
est positif à l'extérieur : $S=\,\big]-\infty\,;\,-\tfrac{1}{2}\big]\cup[2\,;\,+\infty[$
(bornes incluses car l'inégalité est large).
\end{exemple}

\methode{Résoudre une équation se ramenant au second degré}
Repérer le motif répété, poser un changement d'inconnue $X$, résoudre le second
degré en $X$, vérifier les contraintes ($X\ge 0$ pour $X=x^2$), puis revenir à $x$.
\begin{exemple}
$x^4-13x^2+36=0$ : avec $X=x^2$, $X^2-13X+36=0$ donne $X=4$ et $X=9$, d'où
$x\in\{-3\,;\,-2\,;\,2\,;\,3\}$.
\end{exemple}

\methode{Résoudre un système par somme--produit}
Si le système se met sous la forme $\{u+v=S,\ uv=P\}$, écrire l'équation
$t^2-St+P=0$ et résoudre : ses racines sont $u$ et $v$ (à l'ordre près).
\begin{exemple}
$\begin{cases} x+y=10\\ xy=21 \end{cases}$ : $t^2-10t+21=0$, $\Delta=16$, racines
$3$ et $7$. Solutions : $(3\,;\,7)$ et $(7\,;\,3)$.
\end{exemple}
